\chapter{Evaluating \texorpdfstring{$F_{\bm h}$}{Fh}: atoms and the four coherent channels}
\label{chap:unit-cells}

This chapter evaluates the matrix element $F_{\bm h}$ of
Eq.~\eqref{eq:Fh-definition} for the atomic part of the contrast.  The
reference part, which exists only at $\bm h=\bm0$, is evaluated in
Chapter~\ref{chap:specular}.  Both are contributions to the same
$F_{\bm h}$; they are never separate matrix elements.

\section{Where the four channels come from}

Section~\ref{sec:green} showed that the reciprocal exit field entering
Eq.~\eqref{eq:Fh-definition} is $u_-$, the solution of the homogeneous
stratified problem selected by the boundary condition below the stack.  Inside
prepared medium $m$ that solution has two branches, and so does the incident
field:
\begin{equation}
  \bm E_0(z',-\bm k_f)
  =A_{m,f}^{+}\ee^{-\ii k_{z,m,f}z'}+A_{m,f}^{-}\ee^{+\ii k_{z,m,f}z'},
  \qquad
  \bm E_0(z',\bm k_i)
  =A_{m,i}^{+}\ee^{-\ii k_{z,m,i}z'}+A_{m,i}^{-}\ee^{+\ii k_{z,m,i}z'}.
  \label{eq:two-branches}
\end{equation}
Their product in Eq.~\eqref{eq:Fh-definition} therefore contains four terms,
$A_{m,i}^{\sigma_i}A_{m,f}^{\sigma_f}$.  The four channels are not an
additional model: they are the two branches of each reference field, and
nothing more needs to be assumed to justify them.  Let
\begin{equation}
  \sigma_i,\sigma_f\in\{+1,-1\},
  \qquad
  \nu=(\sigma_i,\sigma_f),
\end{equation}
with the fixed channel order
\begin{equation}
  ++,\quad +-,\quad -+,\quad --.
  \label{eq:channel-order}
\end{equation}

\section{Channel wavevectors and coefficients}

For medium $m$, channel $\nu$ has
\begin{equation}
  \bm Q_{m}^{\nu}
  =\left(\Qpar,Q_{z,m}^{\nu}\right),
  \qquad
  Q_{z,m}^{\nu}
  =-\left(\sigma_i k_{z,m,i}+\sigma_f k_{z,m,f}\right),
  \label{eq:channel-qz}
\end{equation}
with $\Qpar=\bm h$ by the Laue condition, Eq.~\eqref{eq:laue}, and
\begin{equation}
  \left(Q_m^\nu\right)^2
  =|\bm h|^2+\left(Q_{z,m}^{\nu}\right)^2.
  \label{eq:channel-q2}
\end{equation}
Neither $Q_{z,m}^{\nu}$ nor $(Q_m^\nu)^2$ is generally real in an absorbing
medium.  Because Eq.~\eqref{eq:renaud-kz} makes $\operatorname{Re}k_{z,m}\le0$,
Eq.~\eqref{eq:channel-qz} gives $Q_z>0$ for the $++$ channel, as it must.

Let $A_{m,i}^{\sigma_i}$ and $A_{m,f}^{\sigma_f}$ be the local Renaud
electric-field amplitudes of Eq.~\eqref{eq:two-branches}, referenced by the
field solver to $z_{m,i}^{\rm ref}$ and $z_{m,f}^{\rm ref}$.  The complete
coefficient multiplying a Fourier factor $\exp(\ii\bm Q_m^\nu\cdot\bm r)$ is
\begin{align}
  \Cchan_{m,\nu}
  ={}&A_{m,i}^{\sigma_i}A_{m,f}^{\sigma_f}
  P_{m,\nu}\nonumber\\
  &\times\exp\!\left\{
    \ii\left[
      \sigma_i k_{z,m,i}z_{m,i}^{\rm ref}
      +\sigma_f k_{z,m,f}z_{m,f}^{\rm ref}
    \right]
  \right\}.
  \label{eq:channel-coefficient}
\end{align}
The reference phase in Eq.~\eqref{eq:channel-coefficient} is not optional: the
stored $A^\pm$ are local amplitudes, not amplitudes referenced to the global
crystallographic origin.

For the right-handed Renaud polarization basis used by orGUI, the bilinear
polarization contractions are
\begin{align}
  P_{ss}^{\sigma_i\sigma_f}
    &=\cos\varphi,\nonumber\\
  P_{sp}^{\sigma_i\sigma_f}
    &=-\sigma_f\sin\alpha_f\sin\varphi,\nonumber\\
  P_{ps}^{\sigma_i\sigma_f}
    &=-\sigma_i\sin\alpha_i\sin\varphi,\nonumber\\
  P_{pp}^{\sigma_i\sigma_f}
    &=\sin\alpha_i\sin\alpha_f
      \left[
        \frac{\kzero^2\cos\alpha_i\cos\alpha_f}
             {k_{z,m,i}k_{z,m,f}}
        -\sigma_i\sigma_f\cos\varphi
      \right].
  \label{eq:polarization-contractions}
\end{align}
These are evaluated at the nominal $\bm k_i$ and $\bm k_f$; by
Sec.~\ref{sec:why-not-farfield} that remains accurate to $\sim L/R$ even where
a far-field treatment of the \emph{phase} would fail badly.

The solver's upward internal tangential coefficient for $p$ polarization has
the opposite sign from the public Renaud amplitude.  Thus
\begin{equation}
  A_{p,m}^{+}=T_{p,m}^{+},
  \qquad
  A_{p,m}^{-}=-T_{p,m}^{-}.
  \label{eq:p-sign}
\end{equation}
Implementations should use the solver's analytic $A^\sigma/k_z$ limits in the
$pp$ expression rather than divide numerically at a critical angle.

\section{One generated UnitCell record}

The surface model is flattened into generated unit-cell records $u$.  Before
flattening, every source \texttt{UnitCell} is split into the same ordered
structural-layer views used to build the optical profile.  A record therefore
contains one such layer cell, a crystal-model weight $w_u$, its coherent domain
transforms $d$, and possibly signed domain occupancies $p_{ud}$.  A one-layer
source cell is already a single record.  Let $A_u$ be the source cell's lateral
area and $A_{\rm ref}$ the crystal-wide reference area of
Eq.~\eqref{eq:lattice-expansion}.  The required area conversion is
\begin{equation}
  s_u=\frac{A_{\rm ref}}{A_u}.
  \label{eq:area-scale}
\end{equation}
Current Film and interface domain transforms preserve lateral area.  If a
future model permits in-plane strain, $A_u$ in Eq.~\eqref{eq:area-scale} must be
replaced by the transformed projected area, including its Jacobian.

After stacking and the coherent-domain transform, atom $a$ of record $u$ and
domain $d$ is at the physical Cartesian position $\bm r_{uad}$.  The atom is
assigned to the unique prepared optical interval
\begin{equation}
  m(u,a,d)\quad\hbox{such that}\quad z_{uad}\in I_{m(u,a,d)}.
  \label{eq:medium-assignment}
\end{equation}
A generated unit cell that crosses an optical boundary must therefore be split
by atom.  Assigning one medium to the entire cell would use the wrong $A^\pm$,
$k_z$, polarization contraction, and complex form-factor argument for some of
its atoms.

For channel $\nu$ in medium $m$, define the atomic contribution
\begin{align}
  \mathcal A_{u,m}^{\nu}
  ={}&s_u w_u\Cchan_{m,\nu}
  \sum_d p_{ud}
  \sum_{\substack{a\in u\\m(u,a,d)=m}}
  o_a f_a\!\left((Q_m^\nu)^2,E\right)
  \exp\!\left[-W_a\!\left(\bm Q_m^\nu\right)\right]
  \exp\!\left(\ii\bm Q_m^\nu\mathbin{\cdot}\bm r_{uad}\right).
  \label{eq:atomic-channel}
\end{align}
Here $o_a$ is the atomic occupancy.  For the currently implemented in-plane
and out-of-plane Debye--Waller parameters,
\begin{equation}
  W_a(\bm Q_m^\nu)
  =\frac{B_{a,\parallel}|\bm h|^2
          +B_{a,\perp}(Q_{z,m}^\nu)^2}{16\pi^2}.
  \label{eq:debye-waller}
\end{equation}
Equations~\eqref{eq:atomic-factor}, \eqref{eq:channel-q2}, and
\eqref{eq:debye-waller} are evaluated at complex arguments without taking an
absolute value or a real part.

\section{The atomic contribution to \texorpdfstring{$F_{\bm h}$}{Fh}}

The atomic part of record $u$'s contribution to Eq.~\eqref{eq:Fh-definition} is
\begin{equation}
  \boxed{
  F_{\bm h,u}^{\rm at}
  =\sum_m\sum_{\nu\in\{++,+-,-+,--\}}
   \mathcal A_{u,m}^{\nu},
  }
  \label{eq:Fh-atomic}
\end{equation}
and for $\bm h\neq\bm 0$ this is the whole of it,
\begin{equation}
  F_{\bm h}=\sum_u F_{\bm h,u}^{\rm at},
  \qquad \bm h\neq\bm 0,
  \label{eq:Fh-nonspecular}
\end{equation}
because Eq.~\eqref{eq:coeff-nonspec} carries no reference term.  The specular
coefficient $F_{\bm 0}$ carries one additional contribution and is completed in
Chapter~\ref{chap:specular}.

The four terms in Eq.~\eqref{eq:Fh-atomic} are amplitudes and must be added
before an intensity is formed.  The incorrect expression
$\sum_\nu|\mathcal A_{u,m}^{\nu}|^2$ discards the standing-wave interference
between incident and exit branches that Eq.~\eqref{eq:two-branches} makes
explicit.

\section{Useful limits}

In a terminal substrate, the boundary condition that defines $u_-$ in
Eq.~\eqref{eq:bc-minus} gives
\begin{equation}
  A_{M,i}^{-}=A_{M,f}^{-}=0,
  \label{eq:terminal-branches}
\end{equation}
so only the $++$ channel survives.  In the high-angle Born limit of a finite
cell in vacuum, $A_i^+=A_f^+=1$ and all reflected branches vanish; the same
reduction occurs.  These are limiting cases of the four-channel expression,
not reasons to omit channels in finite graded layers.

If the cell is wholly contained in one medium, Eq.~\eqref{eq:atomic-channel}
reduces to four ordinary structure-factor evaluations at the four complex
values $Q_{z,m}^\nu$.  If it crosses a boundary, there is no single structure
factor that can be multiplied by one set of channel coefficients; the
atom-medium partition in Eq.~\eqref{eq:atomic-channel} is essential.
